% ==============================================================================
% 1. CODIFICACIÓN, IDIOMA Y FUENTES
% ==============================================================================
\usepackage[T1]{fontenc}
\usepackage[utf8]{inputenc}
\usepackage[spanish, es-tabla]{babel} % es-tabla usa "Tabla" en vez de "Cuadro"
\usepackage{csquotes}

% Tipografía base (Times profesional con soporte matemático completo)
\usepackage{newtxtext}
\usepackage{newtxmath}
\usepackage{microtype}                % Ajuste fino de espaciado y justificación

% ==============================================================================
% 2. GEOMETRÍA Y DISEÑO DE PÁGINA (Normas APA / Universitario)
% ==============================================================================
\usepackage[left=2.54cm, right=2.54cm, top=2.54cm, bottom=2.54cm]{geometry}
\usepackage{setspace}
\onehalfspacing                       % Interlineado 1.5 estándar y limpio

% ==============================================================================
% 3. GRÁFICOS, COLORES Y TABLAS
% ==============================================================================
\usepackage{xcolor}
\definecolor{s21Primary}{RGB}{0, 108, 91}       % Color institucional principal
\definecolor{s21Secondary}{RGB}{0, 128, 128}

\usepackage{graphicx}
\usepackage{float}
\usepackage{pdflscape}
\usepackage{longtable}
\usepackage{multirow}
\usepackage{colortbl}
\usepackage{tablefootnote}
\usepackage{tikz}
\usetikzlibrary{calc}

% ==============================================================================
% 4. MATEMÁTICAS Y HERRAMIENTAS AUXILIARES
% ==============================================================================
\usepackage{amsmath}
\usepackage{amsfonts}
%\usepackage{amssymb}
\usepackage{lipsum}

% ==============================================================================
% 5. DATOS DEL DOCUMENTO (EDITAR SOLO ESTE BLOQUE EN CADA PROYECTO)
% ==============================================================================
\newcommand{\docInstitucion}{Universidad Siglo 21}
\newcommand{\docMateria}{Materia Aquí}
\newcommand{\docTitulo}{Título del Trabajo Práctico}
\newcommand{\docSubtitulo}{Subtítulo o Tema Específico}
\newcommand{\docAutor}{Gabriel Hertzan}
\newcommand{\docLegajo}{VINF00000}
\newcommand{\docDocente}{Docente / Tutor}
\newcommand{\docVersion}{1.0}
\newcommand{\docFecha}{\today}

% Variables retrocompatibles con tus archivos anteriores
\author{\docAutor}
\title{\docMateria}
\newcommand{\Version}{\docVersion}
\newcommand{\titulo}{\docTitulo}
\newcommand{\tema}{\docSubtitulo}

% ==============================================================================
% 6. ENCABEZADOS Y PIES DE PÁGINA (fancyhdr)
% ==============================================================================
\usepackage{fancyhdr}
\setlength{\headheight}{26pt}         % Altura suficiente para dos líneas en cabecera

\fancypagestyle{fancy}{
	\fancyhf{}
	% Cabecera
	\lhead{\footnotesize \textbf{\docTitulo} \\ \textit{\textcolor{s21Primary}{\docSubtitulo}}}
	\rhead{}
	\renewcommand{\headrulewidth}{0.4pt}
	
	% Pie de página
	\lfoot{\footnotesize Versión: \docVersion \\ \docAutor}
	\cfoot{\footnotesize \leftmark}
	\rfoot{\footnotesize Pág. \thepage}
	\renewcommand{\footrulewidth}{0.4pt}
}

% ==============================================================================
% 7. BIBLIOGRAFÍA Y ENLACES (Siempre al final del preámbulo)
% ==============================================================================
\usepackage[natbibapa]{apacite}
\usepackage[hidelinks, pdfencoding=auto]{hyperref}

% Metadatos PDF para visores (Acrobat, navegadores, etc.)
\hypersetup{
	pdftitle={\docTitulo},
	pdfauthor={\docAutor},
	pdfsubject={\docMateria}
}